\section{Two-site and positive-length physical blocking}

The two-site block is the virtual contraction of two copies of $K$:
% Formula: M^{(i_1,i_2),(j_1,j_2)}_{a,c}
% =\sum_b K^{i_1,j_1}_{a,b}K^{i_2,j_2}_{b,c}.
% Ink-to-index: the square M carries the grouped ket word (i_1,i_2), grouped
% bra word (j_1,j_2), and outer virtual indices a,c.  The two K squares expose
% the four constituent physical indices and share the virtual index b.
% Contractions: the single east--west bond between the K tensors is summed;
% every outer virtual leg and every ket/bra leg remains free.
% Boundary signature: semantically both sides have two outer virtual indices,
% one grouped ket word, and one grouped bra word.  The raw drawing has
% (1,1,1,1) on the M side and (1,1,2,2) on the K side because it displays the
% canonical grouping of the two physical letters rather than hiding them.
% Source verdict: exact house-style reproduction of MPDO_K=MM.png from
% Theorem 4.9, lines 851--856.  The source has no blocking enclosure.
\begin{tenkzequation}
    \tnpic[physical=updown, tensor style=box]{%
        \tn[mpo, up={$(i_1,i_2)$}, down={$(j_1,j_2)$}]{M}}
    \qquad$=$\qquad
    \tnpic[physical=updown, tensor style=box]{%
        \tn[mpo, up=$i_1$, down=$j_1$]{K} &
        \tn[mpo, up=$i_2$, down=$j_2$]{K}}
\end{tenkzequation}
The two ket indices and the two bra indices are grouped into the blocked
physical indices, as in
\cite[Theorem~4.9, lines~851--856]{Cirac2016MPDO_arXiv}.

\begin{definition}[Positive-length physical blocking of an MPO tensor]
    \label{def:mpdo_positive_physical_blocking}
    \lean{MPOTensor.blockTensor}
    \lean{MPOTensor.blockTensor_apply}
    \leanok
    For a positive integer $L$, group $L$ adjacent physical sites and denote
    the resulting local tensor by $M^{[L]}$. It has a length-$L$ ket word
    $I=(i_1,\ldots,i_L)$ and a length-$L$ bra word
    $J=(j_1,\ldots,j_L)$ as its physical indices, and
    \begin{align}
        (M^{[L]})^{I,J}
        &=M^{i_1j_1}\cdots M^{i_Lj_L}.
        \notag
    \end{align}
    Thus the ket and bra words are grouped separately. This is the MPO
    version of the physical blocking used for the basis of normal tensors in
    \cite[lines~317--345]{Cirac2016MPDO_arXiv}. It is a local tensor, not the
    closed-chain operator $O_L(M)$ of
    \cite[lines~962--967]{Cirac2016MPDO_arXiv}.
\end{definition}

\begin{theorem}[Doubled-index form of positive MPO blocking]
    \label{thm:mpdo_positive_blocking_doubled_index}
    \lean{MPOTensor.blockedDoubledIndexEquiv}
    \lean{MPOTensor.decodeBlock_blockedDoubledIndexEquiv}
    \lean{MPOTensor.toMPSTensor_blockTensor}
    \lean{MPOTensor.isInjective_toMPSTensor_blockTensor_iff}
    \leanok
    \uses{def:mpdo_positive_physical_blocking, def:blocked_tensor}
    Let $L>0$. Pairing the ket and bra letters site by site gives the
    canonical identification
    $\{0,\ldots,d^L-1\}^2\simeq\{0,\ldots,d^2-1\}^L$.
    Under this identification, the doubled-index MPS tensor of $M^{[L]}$
    is the physical reindexing of the $L$-blocked doubled-index MPS tensor
    of $M$. In particular, one is injective if and only if the other is.
\end{theorem}

\begin{proof}\leanok
    Decode the blocked ket and bra words, pair their letters at each site,
    and compare the two ordered matrix products. Physical reindexing by a
    bijection preserves the span of the tensor matrices.
\end{proof}

\begin{theorem}[Positive blocking commutes with the MPO product]
    \label{thm:mpdo_positive_blocking_product}
    \lean{MPOTensor.blockTensor_mulTensor}
    \leanok
    \uses{def:mpdo_positive_physical_blocking, def:mpdo_operator_product}
    For $L>0$ and two MPO tensors $M,N$ with the same physical dimension,
    let juxtaposition denote the MPO tensor product obtained by contracting
    the intermediate physical index. Then $(MN)^{[L]}=M^{[L]}N^{[L]}$.
\end{theorem}

\begin{proof}\leanok
    Expand the product tensor and sum over its length-$L$ intermediate
    physical word. Reindexing this word by one blocked physical index gives
    the product of the two blocked tensors.
\end{proof}

\begin{theorem}[Closed MPO chains under physical blocking]
    \label{thm:mpdo_closed_chain_physical_blocking}
    \lean{MPOTensor.mpo_blockTensor_eq_reindex}
    \lean{MPOTensor.IsMPDO.blockTensor}
    \leanok
    \uses{def:mpdo_positive_physical_blocking, def:mpo_family, def:mpdo}
    Let $L\in\N$, and let
    \begin{align}
        \Phi_{N,L}\colon
        \{0,\ldots,d^L-1\}^N
        &\longrightarrow \{0,\ldots,d-1\}^{NL}
        \notag
    \end{align}
    concatenate the $N$ blocked words. For blocked configurations
    $\sigma$ and $\tau$,
    \begin{align}
        \rho^{(N)}(M^{[L]})_{\sigma,\tau}
        &=\rho^{(NL)}(M)_{\Phi_{N,L}(\sigma),\Phi_{N,L}(\tau)}.
        \notag
    \end{align}
    The algebraic identity holds for every $L\in\N$. For $L\geq1$, positivity
    of the original nonempty chain of length $NL$ shows that blocking $L$
    adjacent physical sites of an MPDO again gives an MPDO. This is the
    closed-chain identification implicit
    when the one-site tensor and its two-site blocking are placed in vertical
    canonical form in
    \cite[Appendix~C.4, lines~1952--2017]{Cirac2016MPDO_arXiv}.
\end{theorem}

\begin{proof}\leanok
    \uses{lem:mpo_eval_word_append}
    Write $\sigma_k=(\sigma_{k,0},\ldots,\sigma_{k,L-1})$ and
    $\tau_k=(\tau_{k,0},\ldots,\tau_{k,L-1})$. For each $k$,
    $(M^{[L]})^{\sigma_k,\tau_k}
    =\prod_{j=0}^{L-1}M^{\sigma_{k,j},\tau_{k,j}}$.
    Taking both products in increasing index order gives
    \begin{align}
        \prod_{k=0}^{N-1}(M^{[L]})^{\sigma_k,\tau_k}
        &=\prod_{k=0}^{N-1}\prod_{j=0}^{L-1}
            M^{\sigma_{k,j},\tau_{k,j}}
          =\prod_{r=0}^{NL-1}
            M^{(\Phi_{N,L}(\sigma))_r,(\Phi_{N,L}(\tau))_r}.
        \notag
    \end{align}
    Taking the virtual trace gives the equality in the theorem. The map
    $\Phi_{N,L}$ is a bijection, and simultaneous reindexing of the rows and
    columns preserves positive semidefiniteness.
\end{proof}

\begin{definition}[Two-site blocking of an MPO tensor]
    \label{def:mpdo_two_site_blocking}
    \lean{MPOTensor.blockTwo}
    \leanok
    For an MPO tensor $K$, its two-site blocking $K^{[2]}$ has physical
    indices $(i_0,i_1)$ and $(j_0,j_1)$ and matrices
    \begin{align}
        (K^{[2]})^{(i_0,i_1),(j_0,j_1)}
        &=K^{i_0j_0}K^{i_1j_1}.
        \notag
    \end{align}
    This is the blocking used in
    \cite[Theorem~4.9, lines~851--856]{Cirac2016MPDO_arXiv}.
\end{definition}

\begin{theorem}[Concrete two-site blocking as length-two blocking]
    \label{thm:mpdo_block_two_as_positive_blocking}
    \lean{MPOTensor.blockTwo_eq_blockTensor_reindex}
    \lean{MPOTensor.mpo_blockTwo_eq_reindex_blockTensor}
    \lean{MPOTensor.IsMPDO.blockTwo}
    \leanok
    \uses{def:mpdo_positive_physical_blocking,
        def:mpdo_two_site_blocking, def:mpdo}
    Let $\varphi:\{0,\ldots,d^2-1\}\to
    \{0,\ldots,d-1\}^2$ be the canonical bijection. Then
    $(K^{[2]})^{i,j}
    =(K^{[L]})^{\varphi(i),\varphi(j)}\bigm|_{L=2}$.
    Consequently, for every nonempty chain length $N$, the closed MPO of
    $K^{[2]}$ is obtained from the closed MPO of the general length-two block by
    applying $\varphi$ independently at every ket and bra site. In
    particular, if $K$ generates matrix product density operators, then so
    does $K^{[2]}$.
\end{theorem}

\begin{proof}\leanok
    Decode a blocked index as an ordered pair. The two local tensors are
    then the same product $K^{i_0j_0}K^{i_1j_1}$. Applying this equality at
    every site gives the closed-chain reindexing, and simultaneous row and
    column reindexing preserves positive semidefiniteness.
\end{proof}

\begin{theorem}[MPO blocking preserves literal CPSV doubled-index form]
    \label{thm:mpdo_literal_cpsv_positive_blocking}
    \lean{MPOTensor.IsCPSVCanonicalForm_toMPSTensor_blockTensor}
    \lean{MPOTensor.IsCPSVCanonicalForm_toMPSTensor_blockTwo}
    \leanok
    \uses{def:mpdo_positive_physical_blocking, def:mpdo_two_site_blocking,
        thm:mpdo_positive_blocking_doubled_index,
        thm:mpdo_block_two_as_positive_blocking,
        thm:cpsv_physical_reindexing,
        thm:cpsv_literal_canonical_form_positive_blocking}
    Let $A_M$ be the doubled-index MPS tensor associated to an MPO tensor
    $M$. If $A_M$ is in literal CPSV canonical form and $L\geq1$, then the
    doubled-index MPS tensor associated to $M^{[L]}$ is again in literal CPSV
    canonical form. The same conclusion holds for the concrete two-site
    blocking $M^{[2]}$.
\end{theorem}

\begin{proof}\leanok
    \uses{def:mpdo_positive_physical_blocking, def:mpdo_two_site_blocking,
        thm:mpdo_positive_blocking_doubled_index,
        thm:mpdo_block_two_as_positive_blocking,
        thm:cpsv_physical_reindexing,
        thm:cpsv_literal_canonical_form_positive_blocking}
    The doubled-index tensor of $M^{[L]}$ is, under the canonical pairing of
    ket and bra words, a physical relabeling of $A_M^{[L]}$:
    \begin{align}
        A_{M^{[L]}}
        &=\psi_L^*(A_M^{[L]}).
        \label{eq:mpdo_cpsv_block_tensor_reindex}
    \end{align}
    Theorem~\ref{thm:cpsv_literal_canonical_form_positive_blocking} gives
    literal CPSV canonical form for $A_M^{[L]}$, and
    Theorem~\ref{thm:cpsv_physical_reindexing} carries it across
    $\psi_L$. For two-site blocking one takes $L=2$ and uses the concrete
    identification of Theorem~\ref{thm:mpdo_block_two_as_positive_blocking}.
\end{proof}

% This is the vertical-decomposition input to Appendix C.4, not the later
% positive-fusion or Gibbs conclusion.
\begin{theorem}[Literal one-site vertical decomposition]
    \label{thm:mpdo_literal_cpsv_vertical_decomposition}
    \lean{MPSTensor.IsCPSVCanonicalForm.exists_cpsvVerticalDecomposition}
    \leanok
    \uses{def:mpdo, def:cpsv_canonical_form, def:vertical_tensor_mpdo,
        def:bnt_cpsv}
    Let $M$ generate matrix product density operators, and suppose that its
    doubled-index MPS tensor is in literal CPSV canonical form. Then there are
    normal tensors $A_\alpha$, positive multiplicities $r_\alpha$, positive
    numbers $c_{\alpha,q}$, and a coisometry $U$ such that the
    $A_\alpha$ form a CPSV basis of normal tensors for $\widetilde M$ and,
    writing
    \begin{align}
        B_v
        &=\bigoplus_\alpha\bigoplus_{q=0}^{r_\alpha-1}
          c_{\alpha,q}A_\alpha^v,
        \label{eq:mpdo_literal_cpsv_vertical_decomposition_block}
    \end{align}
    one has
    \begin{align}
        UU^\dagger&=\Id,
        \notag\\
        U\widetilde M_vU^\dagger&=B_v,
        \notag\\
        \widetilde M_v&=U^\dagger B_vU.
        \label{eq:mpdo_literal_cpsv_vertical_decomposition}
    \end{align}
\end{theorem}

\begin{proof}\leanok
    \uses{thm:cpsv_literal_vertical_bnt_grouping_with_isometries,
        thm:cpsv_literal_normalized_grouped_sector_maps,
        lem:same_mpv_exact_isometric_decomposition}
    Group the normal vertical corners into phase classes and normalize their
    physical sector maps. The exact original-sector reconstruction gives
    positive-length matrix-product-vector equality with the corresponding
    weighted direct sum. Hence the spectral representatives retain the CPSV
    basis-of-normal-tensors property for $\widetilde M$. Flatten the pair
    $(\alpha,q)$ into one finite index. Orthogonality of the normalized sector
    maps gives $UU^\dagger=\Id$; their intertwining identities give the middle
    equality in~(\ref{eq:mpdo_literal_cpsv_vertical_decomposition}), and their
    exact reconstruction gives the last equality.
\end{proof}

\begin{theorem}[Literal two-site vertical decomposition]
    \label{thm:mpdo_literal_cpsv_vertical_decomposition_block_two}
    \lean{MPSTensor.IsCPSVCanonicalForm.%
        exists_cpsvVerticalDecomposition_blockTwo}
    \leanok
    \uses{def:mpdo_two_site_blocking, def:vertical_tensor_mpdo,
        def:bnt_cpsv}
    Under the hypotheses of
    Theorem~\ref{thm:mpdo_literal_cpsv_vertical_decomposition}, the concrete
    two-site block $M^{[2]}$ has a decomposition of the same form:
    there are normal tensors $\widehat A_\alpha$, $r^{[2]}_\alpha>0$,
    $c^{[2]}_{\alpha,q}>0$, and a coisometry $U^{[2]}$ such that
    \begin{align}
        U^{[2]}(U^{[2]})^\dagger&=\Id,
        \notag\\
        U^{[2]}\widetilde{M^{[2]}}_v(U^{[2]})^\dagger
        &=\bigoplus_{\alpha,q}c^{[2]}_{\alpha,q}
          \widehat A_\alpha^v,
        \notag\\
        \widetilde{M^{[2]}}_v
        &=(U^{[2]})^\dagger
          \left(\bigoplus_{\alpha,q}c^{[2]}_{\alpha,q}
          \widehat A_\alpha^v\right)U^{[2]}.
        \label{eq:mpdo_literal_cpsv_vertical_decomposition_block_two}
    \end{align}
    The tensors $\widehat A_\alpha$ form a CPSV basis of normal tensors for
    $\widetilde{M^{[2]}}$.
\end{theorem}

\begin{proof}\leanok
    \uses{thm:mpdo_literal_cpsv_positive_blocking,
        thm:mpdo_block_two_as_positive_blocking,
        thm:mpdo_literal_cpsv_vertical_decomposition}
    Two-site blocking preserves both literal CPSV canonical form and
    matrix-product-density-operator positivity. Apply
    Theorem~\ref{thm:mpdo_literal_cpsv_vertical_decomposition} to $M^{[2]}$.
\end{proof}

\begin{corollary}[Paired literal vertical decompositions]
    \label{cor:mpdo_literal_cpsv_vertical_decomposition_pair}
    \lean{MPSTensor.IsCPSVCanonicalForm.%
        exists_cpsvVerticalDecomposition_pair}
    \leanok
    \uses{thm:mpdo_literal_cpsv_vertical_decomposition,
        thm:mpdo_literal_cpsv_vertical_decomposition_block_two}
    Under the hypotheses of
    Theorem~\ref{thm:mpdo_literal_cpsv_vertical_decomposition}, there exist
    simultaneously a one-site vertical decomposition of $M$ and a two-site
    vertical decomposition of $M^{[2]}$ of the respective forms stated in
    Theorems~\ref{thm:mpdo_literal_cpsv_vertical_decomposition}
    and~\ref{thm:mpdo_literal_cpsv_vertical_decomposition_block_two}.
\end{corollary}

\begin{proof}\leanok
    \uses{thm:mpdo_literal_cpsv_vertical_decomposition,
        thm:mpdo_literal_cpsv_vertical_decomposition_block_two}
    Choose the one-site and two-site decompositions independently and take
    their ordered pair.
\end{proof}

\begin{lemma}[Two-site blocking preserves normalized BNT-refined horizontal form]
    \label{lem:mpdo_horizontal_cf_block_two}
    \lean{MPOTensor.IsHorizontalCF.blockTwo}
    \leanok
    \uses{def:horizontal_cf_mpdo, def:mpdo_two_site_blocking,
        thm:sector_bnt_cf_block_tensor}
    If an MPO tensor $K$ is in normalized BNT-refined horizontal form, then its
    two-site blocking $K^{[2]}$ is again in that form.
\end{lemma}

\begin{proof}
    \leanok
    Block every normal representative in the horizontal sector decomposition
    by two sites and square each copy weight. Positive blocking preserves
    irreducibility, left-canonical normalization, eventual linear
    independence of the representative states, and inequivalence of distinct
    representatives. Normalized self-overlap is preserved by
    $O_{A_j^{[2]}A_j^{[2]}}(N)=O_{A_jA_j}(2N)\longrightarrow1$.
    Hence the blocked sector decomposition is again a BNT canonical form. If
    $\varphi$ is the canonical bijection between a pair of ket--bra indices
    and a length-two word in the doubled alphabet, then
    $(K^{[2]})^{\mathrm{MPS}}=\varphi^*\!\left(\operatorname{block}_2(K^{\mathrm{MPS}})\right)$.
    Thus $\varphi$ only relabels the physical basis. Finally, if $X$ is the
    original block-diagonal virtual gauge and $S$ its sector tensor, word
    evaluation gives $((K^{[2]})^{\mathrm{MPS}})_i=X(\varphi^*(S^{[2]}))_iX^{-1}$.
    Therefore this blocked decomposition and the same copy gauges witness
    normalized BNT-refined horizontal form for $K^{[2]}$.
\end{proof}

\begin{theorem}[One-site and two-site vertical canonical forms]
    \label{thm:mpdo_vertical_cf_before_after_block_two}
    \lean{MPOTensor.verticalCF_and_blockTwo_of_horizontalCF}
    \leanok
    \uses{thm:vertical_cf_of_horizontal_cf,
        thm:mpdo_block_two_as_positive_blocking,
        lem:mpdo_horizontal_cf_block_two}
    Let $K$ be an MPO tensor in normalized BNT-refined horizontal form that
    generates matrix product density operators. Then both $K$ and $K^{[2]}$
    satisfy the
    vertical canonical-form conditions of
    Definition~\ref{def:vertical_cf_mpdo}.
    This is the initial canonical-form step in
    \cite[Appendix~C.4, lines~1951--1956]{Cirac2016MPDO_arXiv}.
\end{theorem}

\begin{proof}\leanok
    The preceding blocking results show that $K^{[2]}$ is in normalized
    BNT-refined horizontal form and generates matrix product density operators. Apply
    Theorem~\ref{thm:vertical_cf_of_horizontal_cf} first to $K$ and then to
    $K^{[2]}$.
\end{proof}

\begin{definition}[Normalized vertical-sector maps]
    \label{def:mpdo_normalized_vertical_sector_maps}
    \lean{MPOTensor.VerticalSectorAlgebra}
    \lean{MPOTensor.VerticalWeightedSectorSpace}
    \lean{MPOTensor.verticalMultiplicityTrace}
    \lean{MPOTensor.normalizedVerticalSectorEmbedding}
    \lean{MPOTensor.normalizedVerticalSectorEmbedding_apply}
    \lean{MPOTensor.verticalSectorPartialTrace}
    \lean{MPOTensor.verticalSectorPartialTrace_apply}
    \leanok
    Let the vertical sectors be labelled by $\alpha$, with simple matrix
    algebra $M_{d_\alpha}$ and positive diagonal multiplicity matrix
    \begin{align}
        \mu_\alpha
        &=\diag(\mu_{\alpha,0},\ldots,\mu_{\alpha,r_\alpha-1}).
        \notag
    \end{align}
    Write $m_\alpha=\tr(\mu_\alpha)$. The normalized
    embedding and the left partial trace are
    \begin{align}
        R_\mu\!\left(\bigoplus_\alpha X_\alpha\right)
        &=\bigoplus_\alpha
          \frac{\mu_\alpha}{m_\alpha}\otimes X_\alpha,
        \notag\\
        \widetilde R_\mu\!\left(\bigoplus_\alpha Y_\alpha\right)
        &=\bigoplus_\alpha \tr_{\mathrm{mult}}(Y_\alpha).
        \notag
    \end{align}
    On a weighted sector,
    $\widetilde R_\mu(\lambda_\alpha\mu_\alpha\otimes X_\alpha)
    =\lambda_\alpha m_\alpha X_\alpha$.
    Thus the partial trace is the canonical extension to the full block
    matrix space of the inverse map displayed in
    \cite[Appendix~C.4, lines~1957--1971]{Cirac2016MPDO_arXiv}.
\end{definition}

\begin{theorem}[Complete positivity of the normalized vertical-sector maps]
    \label{thm:mpdo_vertical_sector_partial_trace_kraus}
    \lean{MPOTensor.verticalSectorSwapEquiv}
    \lean{MPOTensor.normalizedVerticalSectorEmbedding_isKrausDirectSumMap}
    \lean{MPOTensor.verticalSectorPartialTrace_isKrausDirectSumMap}
    \leanok
    \uses{def:mpdo_normalized_vertical_sector_maps,
        def:kraus_map_between_direct_sums}
    Let $(d_\alpha)_\alpha$ and $(r_\alpha)_\alpha$ be finite families of
    dimensions with $r_\alpha>0$, and suppose that every diagonal entry of
    $\mu_\alpha$ is positive. Then the normalized embedding
    \begin{align}
        R_\mu\colon\bigoplus_\alpha M_{d_\alpha}(\C)
        &\longrightarrow \bigoplus_\alpha M_{r_\alpha d_\alpha}(\C),
        \notag\\
        R_\mu((X_\alpha)_\alpha)
        &=\left(\frac{\mu_\alpha}{m_\alpha}\otimes X_\alpha\right)_\alpha
        \notag
    \end{align}
    is a direct-sum Kraus map. Without any condition on the multiplicities,
    the map
    \begin{align}
        \widetilde R\colon\bigoplus_\alpha M_{r_\alpha d_\alpha}(\C)
        &\longrightarrow \bigoplus_\alpha M_{d_\alpha}(\C),
        \notag\\
        \widetilde R((Y_\alpha)_\alpha)
        &=(\tr_{\mathrm{mult}}(Y_\alpha))_\alpha
        \notag
    \end{align}
    is a direct-sum Kraus map. Its restriction to the weighted sectors is
    the retraction in
    \cite[Appendix~C.4, lines~1961--1970]{Cirac2016MPDO_arXiv}.
\end{theorem}

\begin{proof}\leanok
    \uses{thm:mpdo_equiv_reindex_cptp,
        thm:mpdo_state_preparation_kraus_resolution,
        thm:mpdo_controlled_kraus_map_cptp,
        thm:mpdo_controlled_kraus_map_same_block,
        thm:mpdo_controlled_dependent_partial_trace_same_block,
        thm:mpdo_controlled_kraus_map_off_block,
        thm:mpdo_controlled_dependent_partial_trace_cptp,
        lem:is_kraus_cp_comp}
    Exchange the two coordinates within each sector by
    $w(\alpha,q,i)=(\alpha,i,q)$. Let $R_w$ be the corresponding matrix
    reindexing. For the normalized embedding, set
    $\rho_\alpha=m_\alpha^{-1}\mu_\alpha$.
    The hypotheses give $\rho_\alpha\succeq0$ and
    $\tr(\rho_\alpha)=1$. Let $\mathcal P$ be the
    orthogonally controlled preparation of these density matrices. Thus
    \begin{align}
        [\mathcal P(Z)]_{\langle\alpha,(i,q)\rangle,
            \langle\beta,(j,s)\rangle}
        &=
        \begin{cases}
            Z_{\langle\alpha,i\rangle,\langle\alpha,j\rangle}
            \rho_\alpha(q,s), & \alpha=\beta,\\
            0, & \alpha\ne\beta.
        \end{cases}
        \notag
    \end{align}
    The preparation Kraus operators resolve the identity in each sector, so
    $\mathcal P$ is trace-preserving and completely positive. Reindexing by
    $w^{-1}$ changes $X_\alpha\otimes\rho_\alpha$ into
    $\rho_\alpha\otimes X_\alpha$. Consequently, the full-matrix extension
    satisfies
    \begin{align}
        [\widehat R_\mu(Z)]_{\langle\alpha,(q,i)\rangle,
            \langle\beta,(s,j)\rangle}
        &=
        \begin{cases}
            \rho_\alpha(q,s)
            Z_{\langle\alpha,i\rangle,\langle\alpha,j\rangle},
            & \alpha=\beta,\\
            0, & \alpha\ne\beta,
        \end{cases}
        \notag\\
        \widehat R_\mu
        &=R_{w^{-1}}\circ\mathcal P.
        \notag
    \end{align}
    This proves complete positivity of the normalized embedding.

    Now let $\mathcal C_{\tr}$ be the controlled partial trace over
    the $\C^{r_\alpha}$ factors. For a matrix $Y$ on
    $\bigoplus_\alpha(\C^{r_\alpha}\otimes\C^{d_\alpha})$,
    the full-matrix extension satisfies
    \begin{align}
        [\widehat{\widetilde R}(Y)]_{\langle\alpha,i\rangle,
            \langle\beta,j\rangle}
        &=
        \begin{cases}
            \sum_{q=0}^{r_\alpha-1}
            Y_{\langle\alpha,(q,i)\rangle,\langle\alpha,(q,j)\rangle},
            & \alpha=\beta,\\
            0, & \alpha\ne\beta.
        \end{cases}
        \notag
    \end{align}
    The diagonal case is the ordinary right partial trace after applying
    $R_w$, and the off-diagonal case vanishes under the orthogonal sector
    control. Hence
    $\widehat{\widetilde R}=\mathcal C_{\tr}\circ R_w$.
    Equivalence reindexing and the controlled dependent partial trace are
    trace-preserving completely positive maps, so their composition has a
    Kraus representation.
\end{proof}

\begin{lemma}[Retraction of the normalized vertical-sector embedding]
    \label{lem:mpdo_vertical_sector_retraction}
    \lean{MPOTensor.verticalMultiplicityTrace_ne_zero}
    \lean{MPOTensor.verticalSectorPartialTrace_comp_normalizedVerticalSectorEmbedding}
    \leanok
    \uses{def:mpdo_normalized_vertical_sector_maps}
    Suppose that every multiplicity space is nonzero and every diagonal
    entry of $\mu_\alpha$ is positive. Then $m_\alpha\ne0$ for every
    $\alpha$, and $\widetilde R_\mu\circ R_\mu=\id$.
\end{lemma}

\begin{proof}\leanok
    Positivity of the diagonal entries and nonemptiness of the multiplicity
    space give $m_\alpha=\sum_q\mu_{\alpha,q}>0$. The Kronecker-product
    identity for the left partial trace then gives
    \begin{align}
        \tr_{\mathrm{mult}}
        \left(m_\alpha^{-1}\mu_\alpha\otimes X_\alpha\right)
        &=m_\alpha^{-1}\tr(\mu_\alpha)X_\alpha
          =X_\alpha.
        \notag
    \end{align}
\end{proof}

\begin{definition}[Products of weighted vertical contractions]
    \label{def:mpdo_weighted_vertical_contraction_products}
    \lean{MPOTensor.weightedVerticalBondContraction}
    \lean{MPOTensor.weightedVerticalBondContractionProduct}
    \lean{MPOTensor.WeightedVerticalBondContractionProductSpanTop}
    \leanok
    \uses{def:mpdo_normalized_vertical_sector_maps,
        def:mpdo_vertical_bond_matrix_contraction}
    For vertical BNT tensors $M_\alpha$ with simple matrix algebras
    $M_{d_\alpha}$ and a family of complex scalars $(c_\alpha)_\alpha$, define
    \begin{align}
        V_c(X)
        &=\bigoplus_\alpha c_\alpha M_\alpha(X),
        \notag\\
        M_\alpha(X)
        &=\sum_{a,b}X_{ba}M_{\alpha,ab}.
        \notag
    \end{align}
    For each $L\geq0$, let
    \begin{align}
        C_L(V_c)
        &=\spn_{\C}\left\{V_c(X_1)\cdots V_c(X_L):
            X_1,\ldots,X_L\in M_D\right\}
          \subseteq \prod_\alpha M_{d_\alpha}.
        \notag
    \end{align}
    Write $V_m$ for the specialization obtained by taking
    $c_\alpha=m_\alpha=\tr(\mu_\alpha)$.
\end{definition}

\begin{lemma}[Contraction against a matrix unit]
    \label{lem:mpdo_vertical_bond_contraction_matrix_unit}
    \lean{MPSTensor.contractBondMatrix_single_transpose}
    \leanok
    \uses{def:mpdo_vertical_bond_matrix_contraction}
    Let $M=(M_{ab})_{a,b=0}^{D-1}$ be a matrix-product tensor whose physical
    index is identified with the ordered pair $(a,b)$. Then
    $M(E_{ba})=M_{ab}$.
\end{lemma}

\begin{lemma}[Contraction in the encoded product coordinate]
    \label{lem:mpdo_vertical_bond_contraction_encoded_matrix_unit}
    \lean{MPSTensor.contractBondMatrix_single_mod_div}
    \leanok
    \uses{lem:mpdo_vertical_bond_contraction_matrix_unit}
    If $r\in\{0,\ldots,D^2-1\}$ encodes the ordered pair
    $(\lfloor r/D\rfloor,r\bmod D)$, then
    $M\!\left(E_{(r\bmod D),\lfloor r/D\rfloor}\right)=M_r$.
\end{lemma}

\begin{definition}[Sectorwise scaling equivalence]
    \label{def:mpdo_vertical_sector_scaling_equivalence}
    \lean{MPOTensor.verticalSectorScaleEquiv}
    \leanok
    For scalars $c_\alpha\ne0$, the map
    \begin{align}
        \varphi_c\colon\prod_\alpha M_{d_\alpha}
        &\longrightarrow\prod_\alpha M_{d_\alpha},
        \notag\\
        (X_\alpha)_\alpha
        &\longmapsto(c_\alpha X_\alpha)_\alpha
        \notag
    \end{align}
    is a linear equivalence.
\end{definition}

\begin{lemma}[Weighted contraction of matrix-unit words]
    \label{lem:mpdo_weighted_contraction_matrix_unit_words}
    \lean{MPOTensor.weightedVerticalBondContractionProduct_single}
    \leanok
    \uses{def:mpdo_weighted_vertical_contraction_products,
        lem:mpdo_vertical_bond_contraction_encoded_matrix_unit}
    For a word $w=((a_1,b_1),\ldots,(a_L,b_L))$,
    \begin{align}
        V_c(E_{b_1a_1})\cdots V_c(E_{b_La_L})
        &=\bigoplus_\alpha c_\alpha^L M_\alpha^w.
        \notag
    \end{align}
\end{lemma}

\begin{lemma}[Transfer of simultaneous word spanning]
    \label{lem:mpdo_weighted_contraction_span_transfer}
    \lean{MPOTensor.%
        weightedVerticalBondContractionProductSpanTop_of_wordTupleSpanTop}
    \leanok
    \uses{def:mpdo_vertical_sector_scaling_equivalence,
        lem:mpdo_weighted_contraction_matrix_unit_words}
    Fix $L\geq0$ and suppose that every $c_\alpha$ is nonzero. If the
    simultaneous word tuples $(M_\alpha^w)_\alpha$ of length $L$ span
    $\prod_\alpha M_{d_\alpha}$, then
    $C_L(V_c)=\prod_\alpha M_{d_\alpha}$.
\end{lemma}

\begin{theorem}[Nonzero-weight vertical bond contractions generate the sector algebra]
    \label{thm:mpdo_nonzero_weight_vertical_bond_contractions_generate_sector_algebra}
    \lean{MPOTensor.%
        exists_positive_weightedVerticalBondContractionProductSpanTop_of_bnt}
    \leanok
    \uses{def:bnt_cpsv, def:mpdo_weighted_vertical_contraction_products}
    Suppose that $(M_\alpha)_\alpha$ is a basis of normal tensors and that
    $(c_\alpha)_\alpha$ is any family of nonzero complex scalars. Then there
    is a positive integer $L$ such that
    $C_L(V_c)=\prod_\alpha M_{d_\alpha}$.
\end{theorem}

\begin{proof}\leanok
    \uses{lem:mpdo_weighted_contraction_span_transfer,
        thm:cpsv_bnt_blocked_simultaneous_span}
    Choose $L>0$ for which the simultaneous word tuples
    $(M_\alpha^w)_\alpha$ span $\prod_\alpha M_{d_\alpha}$. Since every
    $c_\alpha$ is nonzero, so is $c_\alpha^L$. Apply
    Lemma~\ref{lem:mpdo_weighted_contraction_span_transfer}.
\end{proof}

\begin{theorem}[Vertical bond contractions generate the sector algebra]
    \label{thm:mpdo_vertical_bond_contractions_generate_sector_algebra}
    \lean{MPOTensor.%
        exists_positive_verticalMultiplicityTraceBondContractionProductSpanTop_of_bnt}
    \leanok
    \uses{def:bnt_cpsv, def:mpdo_weighted_vertical_contraction_products}
    Suppose that $(M_\alpha)_\alpha$ is a basis of normal tensors, every
    multiplicity space is nonzero, and every diagonal entry of $\mu_\alpha$
    is positive. Then there is a positive integer $L$ such that
    $C_L(V_m)=\prod_\alpha M_{d_\alpha}=:\mathcal A_1$.
    This is the spanning assertion used in
    \cite[Appendix~C.4, lines~1980--1990]{Cirac2016MPDO_arXiv}.
\end{theorem}

\begin{proof}\leanok
    \uses{lem:mpdo_vertical_sector_retraction,
        thm:mpdo_nonzero_weight_vertical_bond_contractions_generate_sector_algebra}
    Since every multiplicity space is nonzero and every diagonal entry is
    positive,
    $m_\alpha
    =\sum_{q=0}^{\mult(\alpha)-1}\mu_{\alpha,q}>0$.
    Apply the preceding theorem to the nonzero family
    $c_\alpha=m_\alpha=\tr(\mu_\alpha)$.
\end{proof}

\begin{theorem}[Fixed spanning contraction products determine an endomorphism]
    \label{thm:mpdo_fixed_spanning_vertical_contraction_products}
    \lean{MPOTensor.eq_id_of_weightedVerticalBondContractionProducts_fixed}
    \leanok
    \uses{def:mpdo_weighted_vertical_contraction_products}
    Let $F$ be a linear endomorphism of
    $\prod_\alpha M_{d_\alpha}$. Suppose that $C_L(V_c)$ is the whole
    sector algebra and, for every choice of bond matrices
    $X_1,\ldots,X_L$,
    \begin{align}
        F(V_c(X_1)\cdots V_c(X_L))
        &=V_c(X_1)\cdots V_c(X_L).
        \notag
    \end{align}
    Then $F$ is the identity.
\end{theorem}

\begin{proof}\leanok
    Since the displayed products span the domain, every
    $Y\in\prod_\alpha M_{d_\alpha}$ has the form
    $Y=\sum_i\lambda_iV_c(X_1^{(i)})\cdots V_c(X_L^{(i)})$. Hence
    \begin{align}
        F(Y)
        &=\sum_i\lambda_i
          F\!\left(V_c(X_1^{(i)})\cdots V_c(X_L^{(i)})\right)
          =\sum_i\lambda_i
          V_c(X_1^{(i)})\cdots V_c(X_L^{(i)})
          =Y.
        \notag
    \end{align}
    Thus $F$ is the identity.
\end{proof}

\begin{definition}[Products of fixed points]
    \label{def:mpdo_vertical_fixed_point_products}
    \lean{MPOTensor.fixedPointProductSpan}
    \leanok
    Let $F$ be a linear endomorphism of
    $\mathcal A=\prod_\alpha M_{d_\alpha}$, and write
    $\mathcal F=\{X\in\mathcal A:F(X)=X\}$. For $L\geq0$, define
    \begin{align}
        C_L(\mathcal F)
        &=\spn_{\C}\{X_1\cdots X_L:X_1,\ldots,X_L\in\mathcal F\}.
        \notag
    \end{align}
    This is the product space used in the dimension argument of
    \cite[Appendix~C.4, lines~1980--1993]{Cirac2016MPDO_arXiv}.
\end{definition}

\begin{theorem}[Density blocks of a faithful fixed family bound products]
    \label{thm:mpdo_density_blocks_fixed_point_product_dimension}
    \lean{MPOTensor.fixedPointProductSpan_finrank_le_of_densityBlocks}
    \leanok
    \uses{def:mpdo_vertical_fixed_point_products}
    Let $\mathcal A=\prod_\alpha M_{d_\alpha}$, and let
    $F:\mathcal A\to\mathcal A$ be positive and trace preserving.
    Suppose that the trace adjoint $F^*$ satisfies the Schwarz inequality and
    that $F$ fixes a family $\rho=(\rho_\alpha)_\alpha$ for which every
    $\rho_\alpha$ is positive definite. If $\mathcal F=\Fix(F)$, then, for
    every $L>0$, $\dim C_L(\mathcal F)\leq\dim\mathcal F$.
    % Full-support corollary of
    % \cite[Theorem~6.14 and Equation~(6.63)]{Wolf2012Quantum}; the general
    % theorem allows a complementary zero summand and does not require a
    % positive-definite fixed family.  Used in
    % \cite[Appendix~C.4, lines~1980--1995]{Cirac2016MPDO_arXiv} after such a
    % family is obtained.  See
    % docs/paper-gaps/cpsv16_vertical_sector_invertibility.tex.
\end{theorem}

\begin{proof}\leanok
    \uses{thm:direct_sum_full_matrix_extension_compatibility,
        thm:mean_ergodic_projection_density_block_form}
    Extend $F$ to the full matrix algebra by diagonal compression followed by
    block-diagonal embedding, and choose a numbered basis of the ambient
    space. The extension is positive and trace preserving, its trace adjoint
    satisfies the Schwarz inequality, and the embedded family $\rho$ is a
    positive definite fixed point. The fixed points therefore have the form
    \begin{align}
        U\left(\bigoplus_k\sigma_k\otimes X_k\right)U^*,
        \qquad X_k&\in M_{h_k}(\C),
        \notag
    \end{align}
    where $\tr(\sigma_k)=1$. The parametrization
    $(X_k)_k\mapsto\bigoplus_k\sigma_k\otimes X_k$ is injective, since its
    left partial trace is $(X_k)_k$. Every product of $L$ fixed points lies
    in the range of
    \begin{align}
        (Y_k)_k
        &\longmapsto
        U\left(\bigoplus_k\sigma_k^L\otimes Y_k\right)U^*.
        \notag
    \end{align}
    Thus the span of these products has dimension at most
    $\sum_k h_k^2$, which is the dimension of the fixed-point space of the
    extension. Writing $\mathcal E$ for the fixed-point space of the
    full-matrix extension, the injective block-diagonal embedding and
    diagonal compression give the dimension chain
    \begin{align}
        \dim C_L(\mathcal F)
        &\leq \dim C_L(\mathcal E)
          \leq \dim\mathcal E
          \leq \dim\mathcal F.
        \notag
    \end{align}
\end{proof}

\begin{theorem}[Fixed contractions and the dimension of their fixed-point space]
    \label{thm:mpdo_fixed_contractions_product_dimension}
    \lean{MPOTensor.%
        eq_id_of_weightedVerticalBondContractions_fixed_of_productSpan_finrank_le}
    \leanok
    \uses{def:mpdo_weighted_vertical_contraction_products,
        def:mpdo_vertical_fixed_point_products}
    Let $F$ be a linear endomorphism of
    $\mathcal A=\prod_\alpha M_{d_\alpha}$. Suppose that the following three
    conditions hold, with the second required for every bond matrix $X$:
    \begin{align}
        C_L(V_c)&=\mathcal A,
        \notag\\
        F(V_c(X))&=V_c(X),
        \notag\\
        \dim C_L(\mathcal F)&\leq\dim\mathcal F,
        \qquad
        \mathcal F=\{Y\in\mathcal A:F(Y)=Y\}.
        \notag
    \end{align}
    Then $F=\id_{\mathcal A}$.

    This is the finite-dimensional conclusion in
    \cite[Appendix~C.4, lines~1980--1993]{Cirac2016MPDO_arXiv}. It does not
    assume that the products $V_c(X_1)\cdots V_c(X_L)$ are fixed.
    % The displayed dimension inequality is the remaining consequence of the
    % density-weighted fixed-point theorem for the transported maps.
\end{theorem}

\begin{proof}\leanok
    Since every $V_c(X)$ belongs to $\mathcal F$, every product occurring in
    $C_L(V_c)$ belongs to $C_L(\mathcal F)$. Hence
    \begin{align}
        \mathcal A=C_L(V_c)
        &\subseteq C_L(\mathcal F)\subseteq\mathcal A,
        \notag
    \end{align}
    so $C_L(\mathcal F)=\mathcal A$. Therefore
    \begin{align}
        \dim\mathcal A=\dim C_L(\mathcal F)
        &\leq\dim\mathcal F\leq\dim\mathcal A.
        \notag
    \end{align}
    Thus $\mathcal F=\mathcal A$, and $F$ is the identity.
\end{proof}

\begin{theorem}[Faithful fixed family in a finite product]
    \label{thm:mpdo_direct_sum_faithful_fixed_family}
    \lean{Matrix.IsPositiveDirectSumMap.%
        exists_posDef_fixedFamily_of_fixed_product_span}
    \leanok
    \uses{thm:positive_map_faithful_fixed_point_from_product_span,
        thm:direct_sum_full_matrix_extension_compatibility}
    Let $F$ be a positive trace-preserving endomorphism of
    $\mathcal A=\prod_\alpha M_{d_\alpha}$. Suppose that each $V_j$ is
    fixed by $F$ and, for some $L>0$, the identity belongs to
    \begin{align}
        \spn_{\C}
        \{V_{j_1}\cdots V_{j_L}:j_1,\ldots,j_L\in J\}.
        \notag
    \end{align}
    Then there is a fixed family $\rho=(\rho_\alpha)_\alpha$ such that every
    $\rho_\alpha$ is positive definite.

    This is a local finite-product consequence used to formalize the
    fixed-point argument in
    \cite[Appendix~C.4, lines~1980--1993]{Cirac2016MPDO_arXiv}; CPSV16 applies
    Wolf's density-block description directly rather than stating this
    intermediate result.
\end{theorem}

\begin{proof}\leanok
    Extend $F$ to the full matrix algebra by diagonal compression followed
    by block-diagonal embedding. This extension is positive and trace
    preserving. Block-diagonal embedding preserves the identity and every
    product $V_{j_1}\cdots V_{j_L}$, so
    Theorem~\ref{thm:positive_map_faithful_fixed_point_from_product_span}
    gives a positive-definite fixed point $\widehat\rho$ of the extension.
    Fixed points of the extension are block diagonal. Thus
    \begin{align}
        \widehat\rho
        &=\bigoplus_\alpha\rho_\alpha,
        \notag\\
        F(\rho)
        &=\rho.
        \notag
    \end{align}
    Every diagonal block $\rho_\alpha$ of the positive-definite matrix
    $\widehat\rho$ is positive definite.
\end{proof}

\begin{theorem}[Identity criterion from fixed contractions]
    \label{thm:mpdo_fixed_contractions_trace_adjoint_schwarz_identity}
    \lean{MPOTensor.%
        eq_id_of_weightedVerticalBondContractions_fixed_of_traceAdjointSchwarz}
    \leanok
    \uses{thm:mpdo_direct_sum_faithful_fixed_family,
        thm:mpdo_density_blocks_fixed_point_product_dimension,
        thm:mpdo_fixed_contractions_product_dimension}
    Let $F$ be a positive trace-preserving endomorphism of
    $\mathcal A=\prod_\alpha M_{d_\alpha}$ whose trace adjoint satisfies the
    Schwarz inequality. Suppose that, for some $L>0$,
    $C_L(V_c)=\mathcal A$ and $F(V_c(X))=V_c(X)$ for every bond matrix $X$.
    Then $F=\id_{\mathcal A}$.
    No product $V_c(X_1)\cdots V_c(X_L)$ is assumed to be fixed.

    This is the implication used for each transported square composite in
    \cite[Appendix~C.4, lines~1980--1995]{Cirac2016MPDO_arXiv}.
\end{theorem}

\begin{proof}\leanok
    The fixed contractions and their positive-length product span give a
    positive-definite fixed family $\rho=(\rho_\alpha)_\alpha$. Hence
    Theorem~\ref{thm:mpdo_density_blocks_fixed_point_product_dimension}
    gives
    $\dim C_L(\mathcal F)\leq\dim\mathcal F$, where
    $\mathcal F=\Fix(F)$.
    Since $C_L(V_c)=\mathcal A$ and $V_c(X)\in\mathcal F$, the preceding
    theorem gives $F=\id_{\mathcal A}$.
\end{proof}
